\documentclass[11pt]{article}

\usepackage[a4paper,margin=1in]{geometry}
\usepackage{amsmath,amssymb,amsthm,mathtools}
\usepackage{bm}
\usepackage{mathrsfs}
\usepackage{hyperref}
\usepackage{enumitem}
\usepackage{booktabs}
\usepackage{microtype}

\hypersetup{
  colorlinks=true,
  linkcolor=blue,
  citecolor=blue,
  urlcolor=blue
}

\newtheorem{theorem}{Theorem}
\newtheorem{proposition}[theorem]{Proposition}
\newtheorem{lemma}[theorem]{Lemma}
\newtheorem{corollary}[theorem]{Corollary}
\theoremstyle{definition}
\newtheorem{definition}[theorem]{Definition}
\newtheorem{assumption}[theorem]{Assumption}
\newtheorem{remark}[theorem]{Remark}
\newtheorem{conjecture}[theorem]{Conjecture}

\newcommand{\R}{\mathbb{R}}
\newcommand{\E}{\mathbb{E}}
\newcommand{\Prob}{\mathbb{P}}
\newcommand{\supp}{\mathrm{supp}}
\newcommand{\softmax}{\mathrm{softmax}}
\newcommand{\norm}[1]{\left\lVert #1 \right\rVert}
\newcommand{\ip}[2]{\left\langle #1, #2 \right\rangle}
\newcommand{\mc}[1]{\mathcal{#1}}
\newcommand{\Pone}{\mathcal{P}_1}

\title{Theory Skeleton: Mean-Field Training and Gauge Freedom\\in Gaussian Mixture Density Networks with Natural Parameters}
\author{Working draft}
\date{\today}

\begin{document}
\maketitle

\begin{abstract}
This note extends the softplus natural-parameter program from a single Gaussian
head to a fixed-$K$ Gaussian mixture density network (MDN). Each component is
parameterized by Gaussian natural parameters
$(\eta_{1,k},\eta_{2,k})=(\mu_k/\sigma_k^2,-1/(2\sigma_k^2))$, while the second
coordinate is implemented in practice by a raw scale field passed through a
softplus map so that $\eta_{2,k}<0$ automatically. The first theorem proved
here is a finite-horizon shared-path result for noiseless SGD: in the joint
wide/small-step limit, independently initialized wide MDNs track the same
deterministic labeled measure path $\rho_t$ on every fixed interval $[0,T]$.
The second theorem explains why this still does \emph{not} solve the
model-synthesis problem for plain MDNs: even at identical predictive
distributions, no global permutation of the components may align two mixture
representations, and already at the output-head level a linear interpolation can
incur a strictly positive loss barrier. The correct interpretation is that
deterministic mean-field dynamics suppress run-to-run gauge randomness on fixed
finite horizons, but do not canonically remove mixture non-identifiability.
This sharply separates the $K=1$ and $K>1$ stories: the training-side
mean-field bridge survives, but the synthesis-side canonical-alignment premise
used in the single-Gaussian argument can already fail at $K=2$.
\end{abstract}

\section{Executive Summary}

The MDN extension should be read as two separate messages.

\begin{enumerate}[label=(\roman*)]
  \item \textbf{Training-side message.} For a fixed number of mixture
  components $K$, a two-layer Gaussian MDN with softmax gates, Gaussian natural
  parameters, and softplus raw scale heads still admits a finite-horizon
  mean-field program in the raw particle variables. In the noiseless-SGD
  scaling used in two-layer LMC work, this program can be closed with the same
  training-side structure as in the single-Gaussian case.
  \item \textbf{Synthesis-side message.} The deterministic mean-field path is a
  labeled path. It gives repeatability for a chosen initialization law, but it
  does not produce a canonical component decomposition. Global-permutation model
  synthesis can therefore fail even between predictors that represent exactly
  the same conditional density. This is the first place where the $K=1$ and
  $K>1$ stories diverge.
\end{enumerate}

The intended downstream use is clear. The first message supplies the same kind
of training-side bridge that LMC arguments need in the two-layer scalar case
\cite{FerbachEtAl2024}. The second message explains why plain MDNs remain
harder: finite-width randomness can select different representatives of the same
predictive mixture class, and global permutation is too weak to compare them in
general.

\section{Gaussian MDN with Natural Parameters}

Let $(X,Y)\sim P$ with $X\in\mc{X}\subseteq \R^d$ and $Y\in\R$. Fix a component
count $K\ge 2$. For Gaussian natural parameters define
\[
  T_{\rm G}(y):=(y,y^2),
  \qquad
  A_{\rm G}(\eta_1,\eta_2)
  :=
  -\frac{\eta_1^2}{4\eta_2}
  + \frac12 \log\!\Big(\frac{-\pi}{\eta_2}\Big),
  \qquad \eta_2<0,
\]
and write the component negative log-likelihood as
\[
  \ell_{\rm G}(\eta;y)
  :=
  A_{\rm G}(\eta)-\eta_1 y-\eta_2 y^2.
\]

Let $\varphi:\R^p\times \mc{X}\to \R$ be a scalar hidden feature map. A single particle
is
\[
  \theta=(\beta,a,c,\xi)\in \Theta:=\R^K\times\R^K\times\R^K\times \R^p,
\]
where $\beta=(\beta_1,\dots,\beta_K)$ controls the gate logits,
$a=(a_1,\dots,a_K)$ the first natural coordinates, and
$c=(c_1,\dots,c_K)$ the raw scale coordinates.

Fix $\lambda_{\min}>0$ and
\[
  \psi_{\rm sp}(s):=\lambda_{\min}+\log(1+e^s).
\]
At finite width $N$, define for each component $k$
\begin{align}
  g_{k,N}(x)
  &:=
  \frac1N\sum_{i=1}^N \beta_{i,k}\,\varphi(\xi_i,x),
  \\
  \eta_{1,k,N}(x)
  &:=
  \frac1N\sum_{i=1}^N a_{i,k}\,\varphi(\xi_i,x),
  \\
  s_{k,N}(x)
  &:=
  \frac1N\sum_{i=1}^N c_{i,k}\,\varphi(\xi_i,x),
  \\
  \eta_{2,k,N}(x)
  &:=
  -\psi_{\rm sp}(s_{k,N}(x)).
\end{align}
The gate weights are
\[
  \pi_{k,N}(x)
  :=
  \frac{e^{g_{k,N}(x)}}{\sum_{\ell=1}^K e^{g_{\ell,N}(x)}},
\]
and the predictive conditional density is
\begin{equation}
  p_N(y|x)
  :=
  \sum_{k=1}^K
  \pi_{k,N}(x)\,
  \exp\!\big(
    \eta_{1,k,N}(x)y + \eta_{2,k,N}(x)y^2 - A_{\rm G}(\eta_{k,N}(x))
  \big).
  \label{eq:finite-mdn-density}
\end{equation}

The corresponding mean-field fields for a probability measure
$\rho\in \Pone(\Theta)$ are
\begin{align}
  g_{k,\rho}(x)
  &:=
  \int_\Theta \beta_k \varphi(\xi,x)\,\rho(d\theta),
  \\
  \eta_{1,k,\rho}(x)
  &:=
  \int_\Theta a_k \varphi(\xi,x)\,\rho(d\theta),
  \\
  s_{k,\rho}(x)
  &:=
  \int_\Theta c_k \varphi(\xi,x)\,\rho(d\theta),
  \\
  \eta_{2,k,\rho}(x)
  &:=
  -\psi_{\rm sp}(s_{k,\rho}(x)).
\end{align}
We abbreviate $\eta_{k,\rho}(x):=(\eta_{1,k,\rho}(x),\eta_{2,k,\rho}(x))$ and
\[
  \pi_{k,\rho}(x)
  :=
  \frac{e^{g_{k,\rho}(x)}}{\sum_{\ell=1}^K e^{g_{\ell,\rho}(x)}}.
\]

\begin{proposition}[Automatic admissibility of the second natural coordinates]
\label{prop:mdn-domain}
For every probability measure $\rho\in \Pone(\Theta)$, every finite-width
realization, every component $k$, and every $x\in \mc{X}$,
\[
  \eta_{2,k,\rho}(x)\le -\lambda_{\min}<0,
  \qquad
  \eta_{2,k,N}(x)\le -\lambda_{\min}<0.
\]
\end{proposition}

\begin{proof}
By definition, $\psi_{\rm sp}(s)\ge \lambda_{\min}$ for every $s\in\R$.
\end{proof}

For later use define the Gaussian component derivatives
\begin{align}
  q_1(\eta;y)
  &:=
  \partial_{\eta_1}\ell_{\rm G}(\eta;y)
  =
  -\frac{\eta_1}{2\eta_2}-y,
  \\
  q_2(\eta;y)
  &:=
  \partial_{\eta_2}\ell_{\rm G}(\eta;y)
  =
  \frac{\eta_1^2}{4\eta_2^2}
  -
  \frac{1}{2\eta_2}
  -
  y^2.
\end{align}

\begin{definition}[Responsibilities]
For a measure $\rho\in \Pone(\Theta)$, define the responsibility of component
$k$ by
\begin{equation}
  r_{k,\rho}(x,y)
  :=
  \frac{
    \pi_{k,\rho}(x)\,
    \exp\!\big(\eta_{1,k,\rho}(x)y+\eta_{2,k,\rho}(x)y^2
    -A_{\rm G}(\eta_{k,\rho}(x))\big)
  }{
    \sum_{\ell=1}^K
    \pi_{\ell,\rho}(x)\,
    \exp\!\big(\eta_{1,\ell,\rho}(x)y+\eta_{2,\ell,\rho}(x)y^2
    -A_{\rm G}(\eta_{\ell,\rho}(x))\big)
  }.
  \label{eq:responsibility}
\end{equation}
\end{definition}

\begin{lemma}[MDN differential identities]
\label{lem:mdn-differential-identities}
Let
\[
  \ell_{\rm MDN}(g,\eta;y)
  :=
  -\log\!\left[
    \sum_{k=1}^K
    \pi_k(g)\,
    \exp\!\big(
      \eta_{1,k}y+\eta_{2,k}y^2-A_{\rm G}(\eta_k)
    \big)
  \right],
\]
where $\pi_k(g)=e^{g_k}/\sum_{\ell=1}^K e^{g_\ell}$. Then
\begin{align}
  \partial_{g_k}\ell_{\rm MDN}(g,\eta;y)
  &=
  \pi_k(g)-r_k(g,\eta;y),
  \label{eq:mdn-gate-derivative}
  \\
  \partial_{\eta_{1,k}}\ell_{\rm MDN}(g,\eta;y)
  &=
  r_k(g,\eta;y)\,q_1(\eta_k;y),
  \label{eq:mdn-eta1-derivative}
  \\
  \partial_{\eta_{2,k}}\ell_{\rm MDN}(g,\eta;y)
  &=
  r_k(g,\eta;y)\,q_2(\eta_k;y),
  \label{eq:mdn-eta2-derivative}
\end{align}
where $r_k(g,\eta;y)$ denotes the corresponding responsibility.
\end{lemma}

\begin{proof}
Write
\[
  Z(g,\eta;y)
  :=
  \sum_{\ell=1}^K
  \exp\!\big(
    g_\ell-\log\!\sum_{j=1}^K e^{g_j}
    +\eta_{1,\ell}y+\eta_{2,\ell}y^2-A_{\rm G}(\eta_\ell)
  \big).
\]
Then $\ell_{\rm MDN}=-\log Z$. Differentiating $-\log Z$ with respect to $g_k$
gives
\[
  -\frac{1}{Z}\partial_{g_k}Z
  =
  -\frac{1}{Z}\sum_{\ell=1}^K
  \exp(\cdots)(\delta_{k\ell}-\pi_k)
  =
  \pi_k-r_k.
\]
Differentiating with respect to $\eta_{1,k}$ and $\eta_{2,k}$ gives
\[
  -\frac{1}{Z}\partial_{\eta_{j,k}}Z
  =
  r_k\,\partial_{\eta_j}\ell_{\rm G}(\eta_k;y),
  \qquad j\in\{1,2\},
\]
which is exactly \eqref{eq:mdn-eta1-derivative} and
\eqref{eq:mdn-eta2-derivative}.
\end{proof}

\section{Finite-Horizon SGD Mean-Field Dynamics}

The practical training variables are the raw particle parameters
$(\beta_i,a_i,c_i,\xi_i)$, not the transformed natural coordinates. The MDN
mean-field program must therefore be written directly in the raw parameter
space.

\begin{assumption}[Finite-horizon MDN SGD regime]
\label{ass:mdn-sgd}
Fix a horizon $T<\infty$. Assume:
\begin{enumerate}[label=(M\arabic*)]
  \item the step sizes satisfy
  \[
    s_n=\varepsilon\,\xi(n\varepsilon),
  \]
  where $\xi:[0,\infty)\to(0,\infty)$ is bounded and Lipschitz;
  \item the data stream
  \[
    Z_{n+1}=(X_{n+1},Y_{n+1})
  \]
  is i.i.d.\ with law $P$, and $|Y|\le Y_\ast$ almost surely under $P$;
  \item the feature map $\varphi$ is bounded, and $\nabla_\xi\varphi$ is bounded
  and Lipschitz in $\xi$, uniformly in $x$;
  \item the initialization law $\rho_0$ on $\Theta$ has compact support;
  \item the mixture size $K$ is fixed and finite;
  \item every component uses the softplus scale link
  $\psi_{\rm sp}(s)=\lambda_{\min}+\log(1+e^s)$ with $\lambda_{\min}>0$.
\end{enumerate}
\end{assumption}

Write
\[
  K_\varphi:=\sup_{\xi,x} |\varphi(\xi,x)|,
  \qquad
  K_\nabla:=\sup_{\xi,x}\norm{\nabla_\xi\varphi(\xi,x)},
\]
and let $L_\varphi$ and $L_\nabla$ be Lipschitz constants in $\xi$ for
$\varphi$ and $\nabla_\xi\varphi$.

\begin{proposition}[Raw drift formula for the Gaussian MDN]
\label{prop:mdn-raw-drift}
Assume Assumption~\ref{ass:mdn-sgd}. For a measure $\rho$ define
\begin{align*}
  u_{k,\rho}(x,y)
  &:=
  \pi_{k,\rho}(x)-r_{k,\rho}(x,y),
  \\
  v_{1,k,\rho}(x,y)
  &:=
  r_{k,\rho}(x,y)\,q_1(\eta_{k,\rho}(x);y),
  \\
  v_{2,k,\rho}(x,y)
  &:=
  r_{k,\rho}(x,y)\,q_2(\eta_{k,\rho}(x);y).
\end{align*}
Then the mean-field-scaled population gradient flow associated with the MDN risk
has drift $b=(b_\beta,b_a,b_c,b_\xi)$ given componentwise by
\begin{align}
  b_{\beta_k}(\theta,\rho)
  &:=
  -
  \E\big[u_{k,\rho}(X,Y)\,\varphi(\xi,X)\big],
  \label{eq:mdn-drift-beta}
  \\
  b_{a_k}(\theta,\rho)
  &:=
  -
  \E\big[v_{1,k,\rho}(X,Y)\,\varphi(\xi,X)\big],
  \label{eq:mdn-drift-a}
  \\
  b_{c_k}(\theta,\rho)
  &:=
  \E\big[\psi_{\rm sp}'(s_{k,\rho}(X))\,v_{2,k,\rho}(X,Y)\,\varphi(\xi,X)\big],
  \label{eq:mdn-drift-c}
  \\
  b_\xi(\theta,\rho)
  &:=
  -
  \E\Bigg[
    \sum_{k=1}^K
    \Big(
      \beta_k u_{k,\rho}(X,Y)
      +
      a_k v_{1,k,\rho}(X,Y)
      -
      c_k\psi_{\rm sp}'(s_{k,\rho}(X))v_{2,k,\rho}(X,Y)
    \Big)\nabla_\xi\varphi(\xi,X)
  \Bigg].
  \label{eq:mdn-drift-xi}
\end{align}
\end{proposition}

\begin{proof}
Let
\[
  \mc{R}_N(\theta_1,\dots,\theta_N)
  :=
  \E\big[-\log p_N(Y|X)\big],
\]
where $p_N$ is defined in \eqref{eq:finite-mdn-density}. For particle $i$ and
component $k$,
\[
  \partial_{\theta_i} g_{k,N}(x)
  =
  \frac1N\bigl(e_k\varphi(\xi_i,x),0,0,\beta_{i,k}\nabla_\xi\varphi(\xi_i,x)\bigr),
\]
\[
  \partial_{\theta_i}\eta_{1,k,N}(x)
  =
  \frac1N\bigl(0,e_k\varphi(\xi_i,x),0,a_{i,k}\nabla_\xi\varphi(\xi_i,x)\bigr),
\]
and
\[
  \partial_{\theta_i}\eta_{2,k,N}(x)
  =
  -\frac1N\psi_{\rm sp}'(s_{k,N}(x))
  \bigl(0,0,e_k\varphi(\xi_i,x),c_{i,k}\nabla_\xi\varphi(\xi_i,x)\bigr),
\]
where $e_k$ is the $k$th basis vector of $\R^K$. Applying
Lemma~\ref{lem:mdn-differential-identities} and the chain rule gives
\[
  \nabla_{\theta_i}\mc{R}_N
  =
  \frac1N\bigl(G_{\beta,i},G_{a,i},G_{c,i},G_{\xi,i}\bigr),
\]
where
\begin{align*}
  G_{\beta,i}
  &:=
  \bigl(\E[u_{k,\rho_N}(X,Y)\varphi(\xi_i,X)]\bigr)_{k=1}^K,
  \\
  G_{a,i}
  &:=
  \bigl(\E[v_{1,k,\rho_N}(X,Y)\varphi(\xi_i,X)]\bigr)_{k=1}^K,
  \\
  G_{c,i}
  &:=
  \bigl(
    -\E[\psi_{\rm sp}'(s_{k,\rho_N}(X))v_{2,k,\rho_N}(X,Y)\varphi(\xi_i,X)]
  \bigr)_{k=1}^K,
  \\
  G_{\xi,i}
  &:=
  \E\!\Big[
    \sum_{k=1}^K
    \big(
      \beta_{i,k}u_{k,\rho_N}(X,Y)
      +a_{i,k}v_{1,k,\rho_N}(X,Y)
    \\
    &\hspace{8em}
      -c_{i,k}\psi_{\rm sp}'(s_{k,\rho_N}(X))v_{2,k,\rho_N}(X,Y)
    \big)\nabla_\xi\varphi(\xi_i,X)
  \Big].
\end{align*}
Multiplying by the mean-field factor $N$ in gradient flow yields the stated
drift formulas.
\end{proof}

\begin{equation}
  \partial_t \rho_t
  +
  \nabla_\theta\cdot\bigl(\rho_t\,\xi(t)b(\theta,\rho_t)\bigr)
  =
  0,
  \qquad
  \rho_{t=0}=\rho_0.
  \label{eq:mdn-sgd-pde}
\end{equation}

\begin{proposition}[Time-inhomogeneous finite-horizon confinement for the MDN flow]
\label{prop:mdn-confinement}
Assume Assumption~\ref{ass:mdn-sgd} and let
\[
  \xi_{\max}:=\sup_{t\in[0,T]}\xi(t)<\infty.
\]
Define the initial bounds
\begin{align*}
  B_0
  &:=
  \sup_{\theta\in\supp(\rho_0)}\max_{1\le k\le K}|\beta_k|,
  \\
  A_0
  &:=
  \sup_{\theta\in\supp(\rho_0)}\max_{1\le k\le K}|a_k|,
  \\
  C_0
  &:=
  \sup_{\theta\in\supp(\rho_0)}\max_{1\le k\le K}|c_k|,
  \\
  R_0
  &:=
  \sup_{\theta\in\supp(\rho_0)}\norm{\xi}.
\end{align*}
Let $B_t^{\rm sgd},A_t^{\rm sgd},C_t^{\rm sgd},R_t^{\rm sgd}$ solve
\begin{align}
  \dot B_t^{\rm sgd}
  &=
  \xi_{\max}K_\varphi,
  \qquad
  B_{t=0}^{\rm sgd}=B_0,
  \label{eq:B-envelope-mdn-sgd}
  \\
  \Gamma_{1,t}^{\rm sgd}
  &:=
  \frac{K_\varphi A_t^{\rm sgd}}{2\lambda_{\min}}+Y_\ast,
  \label{eq:Gamma1-mdn-sgd}
  \\
  \dot A_t^{\rm sgd}
  &=
  \xi_{\max}K_\varphi \Gamma_{1,t}^{\rm sgd},
  \qquad
  A_{t=0}^{\rm sgd}=A_0,
  \label{eq:A-envelope-mdn-sgd}
  \\
  \Gamma_{2,t}^{\rm sgd}
  &:=
  \frac{K_\varphi^2 (A_t^{\rm sgd})^2}{4\lambda_{\min}^2}
  +
  \frac{1}{2\lambda_{\min}}
  +
  Y_\ast^2,
  \label{eq:Gamma2-mdn-sgd}
  \\
  \dot C_t^{\rm sgd}
  &=
  \xi_{\max}K_\varphi \Gamma_{2,t}^{\rm sgd},
  \qquad
  C_{t=0}^{\rm sgd}=C_0,
  \label{eq:C-envelope-mdn-sgd}
  \\
  \dot R_t^{\rm sgd}
  &=
  \xi_{\max}K K_\nabla
  \big(
    B_t^{\rm sgd}
    +
    A_t^{\rm sgd}\Gamma_{1,t}^{\rm sgd}
    +
    C_t^{\rm sgd}\Gamma_{2,t}^{\rm sgd}
  \big),
  \qquad
  R_{t=0}^{\rm sgd}=R_0.
  \label{eq:R-envelope-mdn-sgd}
\end{align}
Then these envelopes remain finite on $[0,T]$. Every characteristic of the
time-inhomogeneous mean-field flow
\[
  \dot\theta(t)=\xi(t)b(\theta(t),\rho_t)
\]
started from $\supp(\rho_0)$ satisfies
\[
  \begin{aligned}
    \max_k|\beta_k(t)|&\le B_t^{\rm sgd},
    \qquad
    \max_k|a_k(t)|\le A_t^{\rm sgd},
    \\
    \max_k|c_k(t)|&\le C_t^{\rm sgd},
    \qquad
    \norm{\xi(t)}\le R_t^{\rm sgd},
    \qquad t\in[0,T].
  \end{aligned}
\]
Consequently, with
\[
  \begin{aligned}
    B_t^{\Theta,\rm sgd}
    &:=
    \left\{
      (\beta,a,c,\xi):
      \max_k|\beta_k|\le B_t^{\rm sgd},\;
      \max_k|a_k|\le A_t^{\rm sgd},
    \right.
    \\
    &\qquad\left.
      \max_k|c_k|\le C_t^{\rm sgd},\;
      \norm{\xi}\le R_t^{\rm sgd}
    \right\},
    \qquad t\in[0,T].
  \end{aligned}
\]
every solution remains inside
\[
  B_t^{\Theta,\rm sgd}\subseteq B_T^{\Theta,\rm sgd},
\]
and the corresponding fields satisfy
\[
  |g_{k,\rho_t}(x)|\le K_\varphi B_t^{\rm sgd},
  \qquad
  |\eta_{1,k,\rho_t}(x)|\le K_\varphi A_t^{\rm sgd},
  \qquad
  |s_{k,\rho_t}(x)|\le K_\varphi C_t^{\rm sgd}
\]
for all $k,x,t\in[0,T]$.
\end{proposition}

\begin{proof}
Because $0\le \pi_{k,\rho}\le 1$ and $0\le r_{k,\rho}\le 1$,
\[
  |u_{k,\rho}(x,y)|\le 1.
\]
Also $\eta_{2,k,\rho}(x)\le -\lambda_{\min}$ and
$|\eta_{1,k,\rho}(x)|\le K_\varphi A_t^{\rm sgd}$, hence
\[
  |q_1(\eta_{k,\rho}(x);y)|\le \Gamma_{1,t}^{\rm sgd},
  \qquad
  |q_2(\eta_{k,\rho}(x);y)|\le \Gamma_{2,t}^{\rm sgd}.
\]
Since $0<\psi_{\rm sp}'<1$, the drift formulas
\eqref{eq:mdn-drift-beta}--\eqref{eq:mdn-drift-xi} imply
\[
  |\dot\beta_k(t)|\le \xi_{\max}K_\varphi,
  \qquad
  |\dot a_k(t)|\le \xi_{\max}K_\varphi\Gamma_{1,t}^{\rm sgd},
  \qquad
  |\dot c_k(t)|\le \xi_{\max}K_\varphi\Gamma_{2,t}^{\rm sgd},
\]
and
\[
  \norm{\dot\xi(t)}
  \le
  \xi_{\max}K K_\nabla
  \big(
    B_t^{\rm sgd}
    +
    A_t^{\rm sgd}\Gamma_{1,t}^{\rm sgd}
    +
    C_t^{\rm sgd}\Gamma_{2,t}^{\rm sgd}
  \big).
\]
Comparison for scalar ODEs yields the claimed bounds. The field bounds follow
from the estimates
\[
  |g_{k,\rho_t}(x)|
  \le
  \int |\beta_k|\,|\varphi(\xi,x)|\,\rho_t(d\theta)
  \le
  K_\varphi B_t^{\rm sgd},
\]
and similarly for $\eta_{1,k,\rho_t}$ and $s_{k,\rho_t}$.
\end{proof}

\begin{proposition}[Drift regularity on the finite-horizon confinement set]
\label{prop:mdn-drift-regularity}
Assume Assumption~\ref{ass:mdn-sgd}. There exists a constant
$L_T^{\rm mdn}<\infty$ such that for all
$\theta,\bar\theta\in B_T^{\Theta,\rm sgd}$ and all probability measures
$\rho,\nu$ supported on $B_T^{\Theta,\rm sgd}$,
\begin{equation}
  \norm{b(\theta,\rho)-b(\bar\theta,\nu)}
  \le
  L_T^{\rm mdn}\big(\norm{\theta-\bar\theta}+W_1(\rho,\nu)\big),
  \label{eq:mdn-drift-lipschitz}
\end{equation}
and
\[
  \sup_{\theta\in B_T^{\Theta,\rm sgd},\ \supp(\rho)\subseteq B_T^{\Theta,\rm sgd}}
  \norm{b(\theta,\rho)}
  \le
  L_T^{\rm mdn}.
\]
\end{proposition}

\begin{proof}
On $B_T^{\Theta,\rm sgd}$, the maps
\[
  \theta\mapsto \beta_k\varphi(\xi,x),
  \qquad
  \theta\mapsto a_k\varphi(\xi,x),
  \qquad
  \theta\mapsto c_k\varphi(\xi,x)
\]
are uniformly bounded and Lipschitz in $\theta$, uniformly in $x$ and $k$.
Therefore, for every $x$ and $k$,
\[
  |g_{k,\rho}(x)-g_{k,\nu}(x)|
  +
  |\eta_{1,k,\rho}(x)-\eta_{1,k,\nu}(x)|
  +
  |s_{k,\rho}(x)-s_{k,\nu}(x)|
  \le
  C_T^{(1)} W_1(\rho,\nu)
\]
for some finite constant $C_T^{(1)}$.

By Proposition~\ref{prop:mdn-confinement},
\[
  |g_{k,\rho}(x)|\le G_T^{\rm sgd}:=K_\varphi B_T^{\rm sgd},
  \qquad
  |\eta_{1,k,\rho}(x)|\le M_T^{\rm sgd}:=K_\varphi A_T^{\rm sgd},
  \qquad
  -\Lambda_T^{\rm sgd}\le \eta_{2,k,\rho}(x)\le -\lambda_{\min},
\]
where $\Lambda_T^{\rm sgd}:=\psi_{\rm sp}(K_\varphi C_T^{\rm sgd})$. Hence the
component score
\[
  z_{k,\rho}(x,y)
  :=
  g_{k,\rho}(x)
  +
  \eta_{1,k,\rho}(x)y
  +
  \eta_{2,k,\rho}(x)y^2
  -
  A_{\rm G}(\eta_{k,\rho}(x))
\]
is uniformly bounded and Lipschitz in $(g,\eta)$ on the compact set
\[
  [-G_T^{\rm sgd},G_T^{\rm sgd}]^K
  \times
  \Big(
    [-M_T^{\rm sgd},M_T^{\rm sgd}]
    \times
    [-\Lambda_T^{\rm sgd},-\lambda_{\min}]
  \Big)^K
  \times
  [-Y_\ast,Y_\ast].
\]
Since the softmax map on $\R^K$ is globally Lipschitz, both the gate map
$g\mapsto \pi(g)$ and the responsibility map $z\mapsto \softmax(z)$ are
Lipschitz on this compact set. Therefore
\[
  (\rho,x,y)\mapsto u_{k,\rho}(x,y),
  \qquad
  (\rho,x,y)\mapsto v_{1,k,\rho}(x,y),
  \qquad
  (\rho,x,y)\mapsto v_{2,k,\rho}(x,y)
\]
are uniformly bounded and Lipschitz in $\rho$ on the confinement set. Inserting
these bounds into \eqref{eq:mdn-drift-beta}--\eqref{eq:mdn-drift-xi}, and using
the boundedness and Lipschitz continuity of $\varphi$ and $\nabla_\xi\varphi$,
gives \eqref{eq:mdn-drift-lipschitz}. The uniform bound on $b$ follows from the
same estimates with $\bar\theta=\theta$ and $\rho=\nu$.
\end{proof}

\subsection{SGD mean-field approximation and a shared-law common path}

\begin{definition}[One-sample stochastic drift for the MDN]
\label{def:mdn-sample-drift}
For $z=(x,y)\in\mc{X}\times\R$, define
\[
  \hat b(\theta,\rho;z)
  :=
  \bigl(\hat b_\beta(\theta,\rho;z),\hat b_a(\theta,\rho;z),\hat b_c(\theta,\rho;z),\hat b_\xi(\theta,\rho;z)\bigr)
\]
with componentwise coordinates
\begin{align}
  \hat b_{\beta_k}(\theta,\rho;z)
  &:=
  -
  u_{k,\rho}(x,y)\,\varphi(\xi,x),
  \label{eq:mdn-sample-drift-beta}
  \\
  \hat b_{a_k}(\theta,\rho;z)
  &:=
  -
  v_{1,k,\rho}(x,y)\,\varphi(\xi,x),
  \label{eq:mdn-sample-drift-a}
  \\
  \hat b_{c_k}(\theta,\rho;z)
  &:=
  \psi_{\rm sp}'(s_{k,\rho}(x))\,v_{2,k,\rho}(x,y)\,\varphi(\xi,x),
  \label{eq:mdn-sample-drift-c}
  \\
  \hat b_\xi(\theta,\rho;z)
  &:=
  -
  \sum_{k=1}^K
  \Big(
    \beta_k u_{k,\rho}(x,y)
    +
    a_k v_{1,k,\rho}(x,y)
    -
    c_k\psi_{\rm sp}'(s_{k,\rho}(x))v_{2,k,\rho}(x,y)
  \Big)\nabla_\xi\varphi(\xi,x).
  \label{eq:mdn-sample-drift-xi}
\end{align}
\end{definition}

\begin{proposition}[Unbiased raw stochastic update for the MDN]
\label{prop:mdn-unbiased-sample-drift}
Under Assumption~\ref{ass:mdn-sgd},
\[
  b(\theta,\rho)=\E_Z[\hat b(\theta,\rho;Z)],
\]
where $Z=(X,Y)\sim P$. Consequently the noiseless SGD iteration can be written
as
\begin{equation}
  \theta_i^{n+1}
  =
  \theta_i^n
  +
  s_n\,\hat b(\theta_i^n,\hat\rho_n^N;Z_{n+1}),
  \qquad
  \hat\rho_n^N:=\frac1N\sum_{i=1}^N \delta_{\theta_i^n}.
  \label{eq:mdn-sgd-update}
\end{equation}
\end{proposition}

\begin{proof}
Take expectations of
\eqref{eq:mdn-sample-drift-beta}--\eqref{eq:mdn-sample-drift-xi} with respect
to $Z=(X,Y)\sim P$ and compare with
\eqref{eq:mdn-drift-beta}--\eqref{eq:mdn-drift-xi}.
\end{proof}

\begin{proposition}[Discrete-time finite-horizon confinement for MDN SGD]
\label{prop:mdn-sgd-discrete-confinement}
Assume Assumption~\ref{ass:mdn-sgd}. Define deterministic sequences
$B_n^\varepsilon,A_n^\varepsilon,C_n^\varepsilon,R_n^\varepsilon$ by
\begin{align}
  B_{n+1}^\varepsilon
  &=
  B_n^\varepsilon + s_n K_\varphi,
  \qquad
  B_0^\varepsilon=B_0,
  \label{eq:mdn-discrete-B-envelope}
  \\
  A_{n+1}^\varepsilon
  &=
  A_n^\varepsilon
  +
  s_n K_\varphi
  \left(
    \frac{K_\varphi A_n^\varepsilon}{2\lambda_{\min}}+Y_\ast
  \right),
  \qquad
  A_0^\varepsilon=A_0,
  \label{eq:mdn-discrete-A-envelope}
  \\
  C_{n+1}^\varepsilon
  &=
  C_n^\varepsilon
  +
  s_n K_\varphi
  \left(
    \frac{K_\varphi^2 (A_n^\varepsilon)^2}{4\lambda_{\min}^2}
    +
    \frac{1}{2\lambda_{\min}}
    +
    Y_\ast^2
  \right),
  \qquad
  C_0^\varepsilon=C_0,
  \label{eq:mdn-discrete-C-envelope}
  \\
  R_{n+1}^\varepsilon
  &=
  R_n^\varepsilon
  +
  s_n K K_\nabla
  \Biggl[
    B_n^\varepsilon
    +
    A_n^\varepsilon
    \left(
      \frac{K_\varphi A_n^\varepsilon}{2\lambda_{\min}}+Y_\ast
    \right)
    \nonumber\\
  &\qquad\qquad\qquad
    +
    C_n^\varepsilon
    \left(
      \frac{K_\varphi^2 (A_n^\varepsilon)^2}{4\lambda_{\min}^2}
      +
      \frac{1}{2\lambda_{\min}}
      +
      Y_\ast^2
    \right)
  \Biggr],
  \qquad
  R_0^\varepsilon=R_0.
  \label{eq:mdn-discrete-R-envelope}
\end{align}
Then every SGD iterate of \eqref{eq:mdn-sgd-update} satisfies
\[
  \begin{aligned}
    \max_{1\le k\le K}|\beta_{i,k}^n|&\le B_n^\varepsilon,
    \qquad
    \max_{1\le k\le K}|a_{i,k}^n|\le A_n^\varepsilon,
    \\
    \max_{1\le k\le K}|c_{i,k}^n|&\le C_n^\varepsilon,
    \qquad
    \norm{\xi_i^n}\le R_n^\varepsilon,
    \qquad
    \forall i,\ \forall n\le T/\varepsilon.
  \end{aligned}
\]
Moreover, the discrete envelopes are uniformly bounded by the continuous
majorants from Proposition~\ref{prop:mdn-confinement}; in particular every SGD
iterate remains in $B_T^{\Theta,\rm sgd}$.
\end{proposition}

\begin{proof}
We argue by induction on $n$. Assume
\[
  \max_{1\le k\le K}|\beta_{i,k}^n|\le B_n^\varepsilon,
  \qquad
  \max_{1\le k\le K}|a_{i,k}^n|\le A_n^\varepsilon,
  \qquad
  \max_{1\le k\le K}|c_{i,k}^n|\le C_n^\varepsilon,
  \qquad
  \norm{\xi_i^n}\le R_n^\varepsilon
\]
for all $i$. Then, for every component index $k$,
\[
  |g_{k,\hat\rho_n^N}(x)|
  \le
  \frac1N\sum_{i=1}^N |\beta_{i,k}^n|\,|\varphi(\xi_i^n,x)|
  \le
  K_\varphi B_n^\varepsilon,
\]
\[
  |\eta_{1,k,\hat\rho_n^N}(x)|
  \le
  \frac1N\sum_{i=1}^N |a_{i,k}^n|\,|\varphi(\xi_i^n,x)|
  \le
  K_\varphi A_n^\varepsilon,
\]
and
\[
  |s_{k,\hat\rho_n^N}(x)|
  \le
  \frac1N\sum_{i=1}^N |c_{i,k}^n|\,|\varphi(\xi_i^n,x)|
  \le
  K_\varphi C_n^\varepsilon.
\]
Proposition~\ref{prop:mdn-domain} gives
$\eta_{2,k,\hat\rho_n^N}(x)\le -\lambda_{\min}$ for all $k$. Hence
\[
  |u_{k,\hat\rho_n^N}(x,y)|\le 1,
  \qquad
  |v_{1,k,\hat\rho_n^N}(x,y)|
  \le
  \frac{K_\varphi A_n^\varepsilon}{2\lambda_{\min}}+Y_\ast,
\]
and
\[
  |v_{2,k,\hat\rho_n^N}(x,y)|
  \le
  \frac{K_\varphi^2 (A_n^\varepsilon)^2}{4\lambda_{\min}^2}
  +
  \frac{1}{2\lambda_{\min}}
  +
  Y_\ast^2.
\]
Inserting these bounds into the one-sample update
\eqref{eq:mdn-sgd-update} yields
\[
  \max_{1\le k\le K}|\beta_{i,k}^{n+1}|
  \le
  B_{n+1}^\varepsilon,
  \qquad
  \max_{1\le k\le K}|a_{i,k}^{n+1}|
  \le
  A_{n+1}^\varepsilon,
  \qquad
  \max_{1\le k\le K}|c_{i,k}^{n+1}|
  \le
  C_{n+1}^\varepsilon,
\]
and
\[
  \norm{\xi_i^{n+1}}
  \le
  R_{n+1}^\varepsilon.
\]
This closes the induction.

Since $s_n\le \varepsilon\xi_{\max}$ and the right-hand sides of
\eqref{eq:mdn-discrete-B-envelope}--\eqref{eq:mdn-discrete-R-envelope} are
nondecreasing in the envelopes, the standard Euler-vs-ODE comparison gives
\[
  B_n^\varepsilon\le B_{n\varepsilon}^{\rm sgd},
  \qquad
  A_n^\varepsilon\le A_{n\varepsilon}^{\rm sgd},
  \qquad
  C_n^\varepsilon\le C_{n\varepsilon}^{\rm sgd},
  \qquad
  R_n^\varepsilon\le R_{n\varepsilon}^{\rm sgd}
\]
for all $n\le T/\varepsilon$.
\end{proof}

\begin{proposition}[Uniform stochastic-drift regularity for the MDN]
\label{prop:mdn-sample-drift-regularity}
Assume Assumption~\ref{ass:mdn-sgd}. There exists
$\widehat L_T^{\rm mdn}<\infty$ such that for every sample $z=(x,y)$ in the
support of $P$, all $\theta,\bar\theta\in B_T^{\Theta,\rm sgd}$, and all
probability measures $\rho,\nu$ supported on $B_T^{\Theta,\rm sgd}$,
\begin{align}
  \norm{\hat b(\theta,\rho;z)}
  &\le
  \widehat L_T^{\rm mdn},
  \label{eq:mdn-sample-drift-bound}
  \\
  \norm{\hat b(\theta,\rho;z)-\hat b(\bar\theta,\nu;z)}
  &\le
  \widehat L_T^{\rm mdn}
  \bigl(\norm{\theta-\bar\theta}+W_1(\rho,\nu)\bigr).
  \label{eq:mdn-sample-drift-lipschitz}
\end{align}
\end{proposition}

\begin{proof}
On $B_T^{\Theta,\rm sgd}$, the maps
\[
  \theta\mapsto \beta_k\varphi(\xi,x),
  \qquad
  \theta\mapsto a_k\varphi(\xi,x),
  \qquad
  \theta\mapsto c_k\varphi(\xi,x)
\]
are uniformly bounded and Lipschitz in $\theta$, uniformly in $x$ and $k$.
Hence
\[
  |g_{k,\rho}(x)-g_{k,\nu}(x)|
  +
  |\eta_{1,k,\rho}(x)-\eta_{1,k,\nu}(x)|
  +
  |s_{k,\rho}(x)-s_{k,\nu}(x)|
  \le
  C_T^{(2)}W_1(\rho,\nu)
\]
for some finite constant $C_T^{(2)}$. By Proposition~\ref{prop:mdn-confinement},
all relevant fields lie in the compact set
\[
  G_T^{\rm sgd}:=K_\varphi B_T^{\rm sgd},
  \qquad
  M_T^{\rm sgd}:=K_\varphi A_T^{\rm sgd},
  \qquad
  \Lambda_T^{\rm sgd}:=\psi_{\rm sp}(K_\varphi C_T^{\rm sgd}),
\]
\[
  [-G_T^{\rm sgd},G_T^{\rm sgd}]^K
  \times
  \Big(
    [-M_T^{\rm sgd},M_T^{\rm sgd}]
    \times
    [-\Lambda_T^{\rm sgd},-\lambda_{\min}]
  \Big)^K
  \times
  [-Y_\ast,Y_\ast].
\]
On this set the component score map is bounded and Lipschitz, and the softmax
map is globally Lipschitz. Therefore
\[
  (\rho,x,y)\mapsto u_{k,\rho}(x,y),
  \qquad
  (\rho,x,y)\mapsto v_{1,k,\rho}(x,y),
  \qquad
  (\rho,x,y)\mapsto v_{2,k,\rho}(x,y)
\]
are uniformly bounded and Lipschitz in $\rho$. Inserting these bounds into
\eqref{eq:mdn-sample-drift-beta}--\eqref{eq:mdn-sample-drift-xi}, and using the
boundedness and Lipschitz continuity of $\varphi$ and $\nabla_\xi\varphi$,
gives \eqref{eq:mdn-sample-drift-bound} and
\eqref{eq:mdn-sample-drift-lipschitz}.
\end{proof}

\begin{theorem}[Training-side theorem: finite-horizon SGD mean-field approximation for Gaussian MDNs]
\label{thm:mdn-sgd-mean-field}
Assume Assumption~\ref{ass:mdn-sgd}. Consider the noiseless SGD iterates
\eqref{eq:mdn-sgd-update} and let $\rho_t$ be the deterministic measure-valued
path solving \eqref{eq:mdn-sgd-pde}. Then:
\begin{enumerate}[label=(\roman*)]
  \item the PDE \eqref{eq:mdn-sgd-pde} is well posed on $[0,T]$;
  \item there exists an error term
  $\delta_{N,\varepsilon}^{\rm mdn}(T)\to 0$ as $N\to\infty$ and
  $\varepsilon\downarrow 0$ such that
  \[
    \sup_{n\le T/\varepsilon}
    W_1\!\bigl(\hat\rho_n^N,\rho_{n\varepsilon}\bigr)
    \le
    \delta_{N,\varepsilon}^{\rm mdn}(T)
  \]
  with high probability;
  \item in bounded-support regimes, one obtains the same qualitative structure
  as in the Mei/Ferbach four-dynamics argument: a Monte-Carlo term of order
  $N^{-1/2}$ together with discretization and stochastic-gradient terms of
  order $\varepsilon$ and $\sqrt{\varepsilon}$.
\end{enumerate}
\end{theorem}

\begin{proof}[Proof sketch]
Propositions~\ref{prop:mdn-confinement}, \ref{prop:mdn-drift-regularity},
\ref{prop:mdn-unbiased-sample-drift},
\ref{prop:mdn-sgd-discrete-confinement}, and
\ref{prop:mdn-sample-drift-regularity} verify the model-specific hypotheses
needed for the bounded-support four-dynamics comparison of
Mei--Misiakiewicz--Montanari~\cite{MeiMisiakiewiczMontanari2019}. The same
finite-horizon reduction underlies the SGD-side LMC argument of
Ferbach et al.~\cite{FerbachEtAl2024}. Applying that machinery to the present
MDN drift yields the stated approximation theorem.
\end{proof}

\begin{corollary}[Training-side corollary: shared-law common path for SGD in Gaussian MDNs]
\label{cor:mdn-sgd-common-path}
Assume Assumption~\ref{ass:mdn-sgd}. Let $\hat\rho_{A,n}^N,\hat\rho_{B,n}^N$ be
two independent width-$N$ SGD runs of \eqref{eq:mdn-sgd-update}, both
initialized i.i.d.\ from the same law $\rho_0$ and trained over the same
horizon $[0,T]$. Then
\[
  \sup_{n\le T/\varepsilon}
  W_1\!\bigl(\hat\rho_{A,n}^N,\hat\rho_{B,n}^N\bigr)
  \le
  2\,\delta_{N,\varepsilon}^{\rm mdn}(T)
\]
with high probability. In particular,
\[
  \sup_{n\le T/\varepsilon}
  W_1\!\bigl(\hat\rho_{A,n}^N,\hat\rho_{B,n}^N\bigr)\to 0
\]
in probability as $N\to\infty$ and $\varepsilon\downarrow 0$.
\end{corollary}

\begin{proof}
Apply Theorem~\ref{thm:mdn-sgd-mean-field} to each run and use the triangle
inequality.
\end{proof}

\begin{remark}[How to read the SGD common-path theorem]
Theorem~\ref{thm:mdn-sgd-mean-field} and
Corollary~\ref{cor:mdn-sgd-common-path} should be read as \emph{repeatability in
a chosen labeled gauge}. For a fixed initialization law $\rho_0$ and fixed
finite horizon $[0,T]$, the wide/small-step limit of noiseless SGD is
deterministic. In that precise sense, run-to-run gauge randomness is frozen in
the mean-field limit. What the theorem does \emph{not} say is that the MDN
component decomposition has become canonical.
\end{remark}

\section{Deterministicity Does Not Canonically Fix the Mixture Gauge}

\begin{definition}[Labeled gauge, predictive equivalence, and global permutation]
For a $K$-component MDN written in the raw coordinates
$(\beta,a,c,\xi)$, the \emph{labeled gauge} means that the component labels
$k=1,\dots,K$ are kept fixed and compared coordinatewise. Two $K$-component
MDN predictors are \emph{predictively equivalent} if they induce the same
conditional density $p(y|x)$ for every $(x,y)$. The \emph{mixture gauge}
refers to the resulting non-identifiability inside such a predictive
equivalence class: multiple labeled component decompositions can represent the
same predictor. A \emph{global permutation} acts on a predictor by applying the
same permutation of $\{1,\dots,K\}$ to the component indices for every $x$.
\end{definition}

\begin{remark}[What the common-path theorem fixes and what it does not]
Theorem~\ref{thm:mdn-sgd-mean-field} fixes the dynamics only in the labeled
gauge: for a fixed initialization law $\rho_0$ and finite horizon $[0,T]$, the
wide/small-step limit selects a deterministic labeled path. It does not choose
a canonical representative inside the mixture gauge class of the corresponding
predictor. For $K=1$ this distinction is absent or much smaller; for $K>1$ it
is the new synthesis obstruction.
\end{remark}

\begin{theorem}[Global-permutation linear synthesis can fail already at the head level]
\label{thm:mdn-obstruction}
Fix $\sigma^2>0$ and let $\phi_\sigma(\cdot-m)$ denote the Gaussian density
with mean $m$ and variance $\sigma^2$. On any input domain containing both
negative and positive points, consider the two $K=2$ MDN predictors
\[
  p_A(y|x)
  :=
  \frac12 \phi_\sigma(y-x)+\frac12 \phi_\sigma(y+x),
\]
and
\[
  p_B(y|x)
  :=
  \frac12 \phi_\sigma(y-|x|)+\frac12 \phi_\sigma(y+|x|).
\]
Then:
\begin{enumerate}[label=(\roman*)]
  \item $p_A(y|x)=p_B(y|x)$ for every $(x,y)$, so the two predictors are
  predictively equivalent;
  \item for any global permutation $\tau\in S_2$, the component-parameter
  midpoint between the labeled output-head representations $A$ and $\tau\cdot B$
  produces a predictor $\bar p_\tau$ that differs from the common density
  $p_A=p_B$ on a set of inputs with positive measure;
  \item if the data law gives positive mass to both $\{x<0\}$ and $\{x>0\}$,
  and $Y|X\sim p_A(\cdot|X)$, then the population negative log-likelihood
  satisfies
  \[
    \E[-\log \bar p_\tau(Y|X)]
    >
    \E[-\log p_A(Y|X)]
  \]
  for every global permutation $\tau$.
\end{enumerate}
\end{theorem}

\begin{proof}
Item (i) is immediate: for $x\ge 0$, the two formulas coincide term by term,
while for $x<0$ they coincide after swapping the two components.

For item (ii), write the two labeled representations in natural coordinates.
Both predictors have equal gate weights $\pi_1=\pi_2=1/2$ and constant second
natural coordinates
\[
  \eta_{2,1}=\eta_{2,2}=-\frac{1}{2\sigma^2}.
\]
The first natural coordinates are
\[
  \eta_{1,1}^A(x)=\frac{x}{\sigma^2},
  \qquad
  \eta_{1,2}^A(x)=-\frac{x}{\sigma^2},
\]
and
\[
  \eta_{1,1}^B(x)=\frac{|x|}{\sigma^2},
  \qquad
  \eta_{1,2}^B(x)=-\frac{|x|}{\sigma^2}.
\]
If $\tau$ is the identity, then for every $x<0$ the midpoint satisfies
\[
  \bar\eta_{1,1}(x)
  =
  \frac{1}{2}\left(\frac{x}{\sigma^2}+\frac{-x}{\sigma^2}\right)=0,
  \qquad
  \bar\eta_{1,2}(x)
  =
  \frac{1}{2}\left(\frac{-x}{\sigma^2}+\frac{x}{\sigma^2}\right)=0.
\]
Hence $\bar p_{\mathrm{id}}(y|x)=\phi_\sigma(y)$ for every $x<0$, which differs
from $p_A(y|x)$ whenever $x\ne 0$. If $\tau$ is the transposition, the same
collapse occurs for every $x>0$. Thus every global permutation fails on a set
of inputs with positive measure.

For item (iii), define the midpoint excess risk
\[
  \Delta_\tau
  :=
  \E\big[-\log \bar p_\tau(Y|X)\big]
  -
  \E\big[-\log p_A(Y|X)\big].
\]
Since $Y|X\sim p_A(\cdot|X)$,
\[
  \Delta_\tau
  =
  \E\Big[
    \mathrm{KL}\big(p_A(\cdot|X)\,\|\,\bar p_\tau(\cdot|X)\big)
  \Big].
\]
By item (ii), the conditional densities differ on a set of $x$ values with
positive $P_X$-measure, so the conditional KL divergence is strictly positive
there. Therefore $\Delta_\tau>0$.
\end{proof}

\begin{remark}[Interpretation for future LMC work]
Theorem~\ref{thm:mdn-obstruction} does \emph{not} contradict
Corollary~\ref{cor:mdn-sgd-common-path}. The two results live at different
levels.
\begin{enumerate}[label=(\alph*)]
  \item The finite-horizon common-path theorem says that, for a fixed
  initialization law and fixed labeled parameter space, the wide/small-step
  limit dynamics are deterministic.
  \item The obstruction theorem says that the labeled parameter space itself is
  not canonically quotiented by global permutation. Two predictively equivalent
  representations may still be globally incompatible for linear synthesis.
\end{enumerate}
Thus the natural expectation for plain MDNs is:
\begin{quote}
In the joint wide/small-step mean-field limit and at fixed finite horizon,
run-to-run gauge randomness should be suppressed for a fixed initialization law;
at finite width and with SGD/noise, different microscopic asymmetries can
select different representatives, and the plain global-permutation route to
linear mode connectivity can fail.
\end{quote}
For a single Gaussian head, the corresponding training-side common-path
statement is much closer to what the synthesis argument needs, because there is
no analogous nontrivial mixture gauge to quotient. For plain MDNs, this is
exactly the step where a naive extension of the $K=1$ story breaks.
\end{remark}

\section{Toward a Gauge-Dominance Theorem}

The strongest version of the MDN-LMC story is not merely that gauge directions
\emph{exist}. It is that, once the ordinary mean-field approximation error has
been controlled, the \emph{remaining} failure of parameter-space synthesis is
dominated by those gauge directions. This section records the theorem template
needed to make that statement rigorous.

The terminology here is structural. The gauge is not created by finite width:
the mixture non-identifiability already exists at the level of the predictor
class. Finite width and SGD noise matter because different runs can select
different representatives of the same predictive class. Accordingly,
$\Delta_{N,\mathrm{g}}$ should be read as the projection of the finite-width
run-to-run discrepancy onto the low-curvature gauge directions, while
$\Delta_{N,\mathrm{st}}$ is the transverse finite-width error.

\begin{definition}[Finite-width synthesis barrier]
Let $\Theta_A^N,\Theta_B^N$ be two trained width-$N$ particle systems after a
chosen neuron alignment procedure. Their linear-interpolation barrier is
\[
  \mathrm{Bar}_N(\Theta_A^N,\Theta_B^N)
  :=
  \sup_{\alpha\in[0,1]}
  \mc{R}\big((1-\alpha)\Theta_A^N+\alpha \Theta_B^N\big)
  -
  \max\{\mc{R}(\Theta_A^N),\mc{R}(\Theta_B^N)\},
\]
where $\mc{R}$ is the population MDN risk.
\end{definition}

\begin{assumption}[Stable/gauge decomposition at the reference path]
\label{ass:gauge-dominance}
Fix a terminal time $T<\infty$ and let $\rho_T$ be the deterministic mean-field
law from Theorem~\ref{thm:mdn-sgd-mean-field}. Assume:
\begin{enumerate}[label=(GD\arabic*)]
  \item there is a local decomposition of perturbations around $\rho_T$ into
  stable directions and gauge directions,
  \[
    \Delta = \Delta_{\mathrm{st}}+\Delta_{\mathrm{g}},
    \qquad
    \Delta_{\mathrm{st}}\in E_{\mathrm{st}},
    \quad
    \Delta_{\mathrm{g}}\in E_{\mathrm{g}};
  \]
  \item on $E_{\mathrm{st}}$, the second variation of the population risk is
  coercive:
  \[
    \delta^2 \mc{R}(\rho_T)[h,h]
    \ge
    m_T \norm{h}_{\mathrm{st}}^2
    \qquad \forall h\in E_{\mathrm{st}}
  \]
  for some $m_T>0$;
  \item on $E_{\mathrm{g}}$, the predictive map is flat to first order and the
  risk curvature is much smaller:
  \[
    D\Pi(\rho_T)[h]=0
    \qquad\text{and}\qquad
    \delta^2 \mc{R}(\rho_T)[h,h]
    \le
    \kappa_T \norm{h}_{\mathrm{g}}^2
    \qquad \forall h\in E_{\mathrm{g}},
  \]
  with $\kappa_T\ll m_T$;
  \item after alignment, the two-run finite-width discrepancy satisfies
  \[
    \Delta_N
    :=
    \Theta_A^N-\Theta_B^N
    =
    \Delta_{N,\mathrm{st}}+\Delta_{N,\mathrm{g}},
  \]
  where
  \[
    \norm{\Delta_{N,\mathrm{st}}}=O_{\Prob}(\varepsilon_N),
    \qquad
    \norm{\Delta_{N,\mathrm{g}}}=O_{\Prob}(\gamma_N),
    \qquad
    \varepsilon_N=o(\gamma_N);
  \]
  \item the barrier admits a local lower bound of the form
  \[
    \mathrm{Bar}_N
    \ge
    c_T \norm{\Delta_{N,\mathrm{g}}}^2
    -
    C_T \norm{\Delta_{N,\mathrm{st}}}^2
    +
    o_{\Prob}\!\big(\norm{\Delta_N}^2\big)
  \]
  for some constants $c_T,C_T>0$.
\end{enumerate}
\end{assumption}

\begin{conjecture}[Gauge-dominant finite-width failure for plain MDNs]
\label{conj:gauge-dominance}
Assume the setting of Corollary~\ref{cor:mdn-sgd-common-path} together with
Assumption~\ref{ass:gauge-dominance}. Then the finite-width synthesis barrier
for plain MDNs is asymptotically governed by the gauge component:
\[
  \mathrm{Bar}_N
  =
  c_T \norm{\Delta_{N,\mathrm{g}}}^2
  +
  o_{\Prob}\!\big(\norm{\Delta_{N,\mathrm{g}}}^2\big).
\]
In particular, if the gauge block is nondegenerate in probability, then there
exist constants $\delta,p_0>0$ such that
\[
  \liminf_{N\to\infty}
  \Prob\!\left(
    \mathrm{Bar}_N \ge \delta\,\gamma_N^2
  \right)
  \ge
  p_0.
\]
Thus, once the ordinary mean-field tracking error has been separated out, the
observed run-to-run failure of linear synthesis is explained by gauge
fluctuations rather than by the transverse finite-width error.
\end{conjecture}

\begin{remark}[Why this distinguishes $K=1$ from $K>1$]
The conjecture above is genuinely about the mixture case. For a single Gaussian
head, the nontrivial gauge space is absent or much smaller, so the
$\Delta_{N,\mathrm{g}}$ term disappears and the barrier is controlled by the
ordinary finite-width error alone. The MDN difficulty is therefore not merely
that finite-width approximation is imperfect in both cases; it is that
$K>1$ introduces additional low-curvature directions that can dominate the
observed synthesis failure. In particular, the training-side common-path
argument from the $K=1$ paper can plausibly be extended to $K>1$ in the
wide-limit finite-horizon regime, but the synthesis-side premise used there can
already fail for $K=2$ because no canonical global alignment need exist.
\end{remark}

\begin{remark}[What remains to be proved]
To turn Conjecture~\ref{conj:gauge-dominance} into a theorem, three missing
ingredients have to be supplied.
\begin{enumerate}[label=(\alph*)]
  \item A local quotient geometry for plain MDNs, identifying the stable block
  $E_{\mathrm{st}}$ and the gauge block $E_{\mathrm{g}}$ near the deterministic
  path.
  \item A finite-width fluctuation theorem, after alignment, that shows the
  stable block is smaller than the gauge block in the regime observed
  experimentally.
  \item A lower bound proving that the interpolation barrier is sensitive to the
  gauge block exactly in the sense postulated in (GD5).
\end{enumerate}
These are not generic mean-field tasks. They are precisely the new ingredients
needed to explain why the mixture case fails while the single-Gaussian case does
not.
\end{remark}

\section{Implications for the Research Program}

The present note supports the following research direction.

\begin{enumerate}[label=(\roman*)]
  \item The \emph{training-side} premise used in the $K=1$ LMC argument still
  survives for plain MDNs. In the current note it is verified in the same
  finite-horizon noiseless-SGD scaling, but only as a deterministic statement
  in a chosen labeled gauge.
  \item The direct extension from $K=1$ to $K>1$ fails at the
  \emph{synthesis-side} premise, not at the mean-field premise. Plain MDNs have
  more symmetry than a single Gaussian head, and global permutation is not rich
  enough to identify all predictively equivalent decompositions.
  \item Consequently, an eventual $K>1$ analogue of the $K=1$ theorem must
  explain why the missing quotient issue is harmless, or else avoid it by
  changing the model class. The next MDN-LMC paper can therefore naturally aim
  for one of two outcomes:
  \begin{enumerate}[label=(\alph*)]
    \item a negative result showing that plain-MDN global-permutation synthesis
    is impossible in general;
    \item a positive result for a restricted synthesis-friendly class, such as
    anchored, hierarchical, or otherwise gauge-breaking MDNs.
  \end{enumerate}
\end{enumerate}

The theorem proved here favors a crisp narrative. For $K=1$, one can hope to
run the full chain from finite-horizon common path to aligned interpolation to
LMC. For $K>1$, the wide-limit common path still looks plausible, but the chain
can break before interpolation because the relevant quotient of MDN parameter
space is no longer captured by global permutation. The mean-field part is not
the hard part anymore; the hard part is identifying the right quotient for
comparing and synthesizing trained mixture models.

\begin{thebibliography}{99}

\bibitem{Bishop2006}
Christopher M. Bishop.
\newblock \emph{Pattern Recognition and Machine Learning}.
\newblock Springer, 2006.

\bibitem{FerbachEtAl2024}
Damien Ferbach, Baptiste Goujaud, Gauthier Gidel, and Aymeric Dieuleveut.
\newblock Proving linear mode connectivity of neural networks via optimal
transport.
\newblock arXiv:2310.19103, 2024.

\bibitem{MeiMisiakiewiczMontanari2019}
Song Mei, Theodor Misiakiewicz, and Andrea Montanari.
\newblock Mean-field theory of two-layers neural networks: Dimension-free
bounds and kernel limit.
\newblock In \emph{Proceedings of the Thirty-Second Conference on Learning
Theory}, pages 2388--2464, 2019.

\bibitem{SirignanoSpiliopoulos2020}
Justin Sirignano and Konstantinos Spiliopoulos.
\newblock Mean field analysis of neural networks: A law of large numbers.
\newblock \emph{SIAM Journal on Applied Mathematics}, 80(2):725--752, 2020.

\bibitem{vanHavreEtAl2015}
Zo\'e van Havre, Nial Friel, Arthur Rousseau, and Chris P. Robert.
\newblock Overfitting Bayesian mixture models with an unknown number of
components.
\newblock \emph{PLoS ONE}, 10(7):e0131739, 2015.

\end{thebibliography}

\end{document}
